\section{Related work}

Prior work on tourism and location-based personalization has largely focused on point-of-interest (POI) recommendation, where systems predict or rank places that match a user profile and context. Surveys of travel and POI recommenders highlight the importance of heterogeneous signals such as mobility patterns, time, geography, and user attributes, as well as evaluation challenges that arise from contextual variability and data sparsity \cite{renjith_extensive_2020,wang_survey_2025}. Recent work in the UMAP community also emphasizes that user preferences can differ before and after a visit, which complicates modeling and suggests that preference signals and content expectations are not static \cite{hofschen_expected_2023}.

With the rise of large language models (LLMs), research has expanded to interactive and generative paradigms for personalization, including LLM-based pipelines for POI recommendation and user profile generation \cite{wang_embracing_2024,wang_point--interest_2025,Li2024SIGIR_LLMNextPOI,noauthor_genup_nodate}. Tourism research has also begun evaluating generative AI for content creation and trip planning \cite{li_generative_2025,zhang_co-creating_2024,hsu_fine-tuned_2024}. However, these works often study destination promotion or end-to-end planning rather than the micro-content users read when deciding whether a POI is relevant. This study addresses this gap by adapting POI descriptions to user interest profiles while keeping core facts stable, measuring perceptions of significance, trustworthiness, and clarity.

